\documentclass[a4paper,12pt]{report}

% --- PACCHETTI ---
\usepackage[utf8]{inputenc}
\usepackage[italian]{babel}
\usepackage{graphicx}     % Per le immagini
\usepackage{hyperref}     % Per i link
\usepackage{geometry}     % Per i margini
\usepackage{listings}     % Per il codice
\usepackage{xcolor}       % Per i colori del codice
\usepackage{float}        % Per posizionare le tabelle/figure
\usepackage{titlesec}     % Per lo stile dei titoli
\usepackage{booktabs}     % Per tabelle più belle
\usepackage{fancyhdr}     % Per intestazioni e piè di pagina

% --- PACCHETTI PER DIAGRAMMI (TIKZ) ---
\usepackage{tikz}
\usetikzlibrary{shapes.geometric, arrows, positioning, calc}

% --- CONFIGURAZIONE GEOMETRIA (RISOLVE PAGINA VUOTA) ---
% Impostiamo solo i margini e l'altezza dell'header, lasciando a geometry il calcolo del resto.
\geometry{
 a4paper,
 margin=2.5cm,      % Margine uniforme di 2.5cm
 headheight=15pt    % Altezza header sufficiente per fancyhdr
}

% --- RISOLUZIONE SPAZI GRANDI VERTICALI ---
% Evita che LaTeX stiri gli spazi verticali per riempire la pagina a tutti i costi.
\raggedbottom 

% --- STILE CODICE JAVA ---
\definecolor{codegreen}{rgb}{0,0.6,0}
\definecolor{codegray}{rgb}{0.5,0.5,0.5}
\definecolor{codepurple}{rgb}{0.58,0,0.82}
\definecolor{backcolour}{rgb}{0.96,0.96,0.96}

\lstdefinestyle{javastyle}{
    backgroundcolor=\color{backcolour},   
    commentstyle=\color{codegreen},
    keywordstyle=\color{blue},
    numberstyle=\tiny\color{codegray},
    stringstyle=\color{codepurple},
    basicstyle=\ttfamily\footnotesize,
    breakatwhitespace=false,         
    breaklines=true,                 
    captionpos=b,                    
    keepspaces=true,                 
    numbers=left,                    
    numbersep=5pt,                  
    showspaces=false,                
    showstringspaces=false,
    showtabs=false,                  
    tabsize=2,
    frame=single
}
\lstset{style=javastyle}

% --- INTESTAZIONI ---
\pagestyle{fancy}
\fancyhf{}
\rhead{Progetto Helios}
\lhead{Sistemi Distribuiti}
\cfoot{\thepage}

% --- DATI PROGETTO ---
\title{
    \vspace{1cm} % Ridotto leggermente per sicurezza
    \Huge \textbf{HELIOS} \\
    \Large Sistema Distribuito per il Monitoraggio e l'Analisi di Impianti Fotovoltaici Industriali \\
    \vspace{3cm}
   \begin{figure}[h]
    \centering
    \includegraphics[width=0.2\textwidth]{university_logo.png}
    \label{fig:esempio}
\end{figure}
    
    \vspace{3cm}
}
\author{
    \textbf{Studente:} [Simone Battiato] \\
    \textbf{Matricola:} [1000031017] \\
    \vspace{0.5cm} \\
    \textit{Corso di Ingegneria dei Sistemi Distribuiti}
}
\date{Anno Accademico 2025/2026}

\begin{document}

\maketitle

% --- ABSTRACT ---
\begin{abstract}
Il progetto descrive l'architettura, le scelte progettuali e l'implementazione di \textbf{Helios}, un sistema software distribuito per la gestione intelligente di impianti di produzione energetica rinnovabile. Il sistema simula uno scenario reale di monitoraggio per una Comunità Energetica (CER) o un impianto industriale, gestendo volumi di dati significativi (oltre 225 MWh annui).

L'architettura segue il modello Client-Server RESTful, implementando pattern avanzati quali \textit{Remote Facade}, \textit{Remote Proxy}, \textit{DTO}, \textit{Session State} e \textit{Request Batch}. Particolare attenzione è stata posta alla sicurezza (tramite protocollo JWT) e all'integrazione di sistemi di Intelligenza Artificiale Generativa (LLM) per il supporto decisionale.
\end{abstract}

\tableofcontents
\newpage

% -------------------------------------------------------------------
% CAPITOLO 1
% -------------------------------------------------------------------
\chapter{Introduzione e Contesto}

\section{Scenario di Riferimento}
Il progetto nasce per rispondere alle esigenze di un'azienda fittizia operante nel settore logistico che possiede un impianto fotovoltaico da circa 200 kWp installato sulla copertura del proprio stabilimento.

L'azienda necessita di trasformare i dati grezzi provenienti dai sensori (SCADA/IoT) in informazioni strategiche per:
\begin{itemize}
    \item \textbf{Monitoraggio Tecnico:} Rilevare cali di efficienza dovuti a sporcizia o guasti.
    \item \textbf{Controllo Economico:} Calcolare il ROI (Return on Investment) e il risparmio in bolletta.
    \item \textbf{Sostenibilità:} Quantificare le emissioni di $CO_2$ evitate.
\end{itemize}

\section{Obiettivi del Sistema}
Helios si propone di risolvere le criticità legate alla gestione passiva dell'impianto. Gli obiettivi principali sono:
\begin{enumerate}
    \item \textbf{Centralizzazione:} Aggregare dati eterogenei (produzione elettrica, dati meteo) in un unico database.
    \item \textbf{Sicurezza:} Garantire l'accesso ai dati sensibili solo a personale autorizzato tramite token crittografici.
    \item \textbf{Intelligence:} Fornire analisi predittive e diagnostiche utilizzando algoritmi statistici e IA Generativa.
\end{enumerate}

% -------------------------------------------------------------------
% CAPITOLO 2
% -------------------------------------------------------------------
\chapter{Architettura del Sistema}

\section{Panoramica Architetturale}
Il sistema è basato su un'architettura distribuita a due livelli (Two-Tier), progettata per garantire disaccoppiamento e scalabilità.

\begin{figure}[H]
    \centering
    % --- DIAGRAMMA ARCHITETTURALE (TIKZ) ---
    \begin{tikzpicture}[node distance=2cm, auto,
        thick,
        % Stili dei nodi
        component/.style={rectangle, draw=black, fill=blue!10, align=center, rounded corners, minimum height=3em, minimum width=6em},
        database/.style={cylinder, draw=black, shape border rotate=90, aspect=0.25, fill=yellow!10, align=center, minimum height=4em, minimum width=3em},
        external/.style={ellipse, draw=black, fill=red!10, align=center, minimum height=3em, minimum width=5em},
        internet/.style={cloud, draw,cloud puffs=10,cloud puff arc=120, aspect=2, inner ysep=1em, fill=gray!10}
    ]

    % Posizionamento Nodi
    \node [component] (client) {\textbf{Client CLI}\\(Java)};
    \node [component, right=3cm of client] (server) {\textbf{Server Backend}\\(Spring Boot)};
    \node [database, below=1.5cm of server] (db) {PostgreSQL\\(Docker)};
    \node [external, right=2.5cm of server] (llm) {Groq API\\(Cloud LLM)};

    % Connessioni
    \draw [->, thick] ([yshift=2mm]client.east) -- node [above, font=\footnotesize] {HTTP REST} ([yshift=2mm]server.west);
    \draw [<-, thick] ([yshift=-2mm]client.east) -- node [below, font=\footnotesize] {JSON} ([yshift=-2mm]server.west);
    \draw [<->, thick] (server) -- node [right, font=\footnotesize] {JPA} (db);
    \draw [<->, thick] (server) -- node [above, font=\footnotesize] {HTTP} (llm);

    % Etichette extra
    \node[below=0.2cm of client, font=\scriptsize, color=gray] {Presentation Layer};
    \node[above=0.2cm of server, font=\scriptsize, color=gray] {Business Layer};

    \end{tikzpicture}
    \caption{Schema ad alto livello dell'interazione tra i componenti del sistema Helios.}
\end{figure}

\subsection{Lato Server (Backend)}
Il cuore del sistema è sviluppato in \textbf{Java} utilizzando il framework \textbf{Spring Boot}. Espone una serie di API RESTful che fungono da interfaccia per qualsiasi client autorizzato.
\begin{itemize}
    \item \textbf{Tecnologie:} Java 21, Spring Boot 3.5.8, Spring Security, JPA/Hibernate, Springdoc OpenAPI, Spring Cache.
    \item \textbf{Database:} PostgreSQL 16 (containerizzato via Docker) per la persistenza dei dati energetici.
    \item \textbf{Integrazione AI:} Modulo di comunicazione HTTPS con \textbf{Groq API} (\texttt{llama-3.3-70b-versatile}) per l'analisi del linguaggio naturale. Fallback automatico su modelli alternativi in caso di rate limiting.
\end{itemize}

\subsubsection{Componenti Backend}

\paragraph{AuthController}
Gestisce l'autenticazione e l'emissione dei token JWT. È l'unico endpoint pubblico del sistema.
\begin{itemize}
    \item \textbf{Endpoint:} \texttt{POST /api/auth/login}
    \item \textbf{Responsabilità:} Verifica le credenziali (username/password), genere un token JWT con scadenza 24 ore, lo restituisce al client.
    \item \textbf{Logica:} Estratto dalle credenziali il subject (utente), il token viene firmato con algoritmo HMAC-SHA256 ed include una data di scadenza.
\end{itemize}

\paragraph{EnergyController}
Strato REST puro: espone gli endpoint HTTP e delega l'intera elaborazione a \texttt{EnergyService}. Non contiene logica di business.
\begin{itemize}
    \item \textbf{Responsabilità:} Ricevere richieste HTTP autenticate, estrarre i parametri e restituire la risposta di \texttt{EnergyService}.
    \item \textbf{Endpoints:} \texttt{/daily-report-session}, \texttt{/batch-suggestions}, \texttt{/wind-impact}, \texttt{/forecast}, \texttt{/financial-report}, \texttt{/peak-hours}, \texttt{/smart-analysis}.
\end{itemize}

\paragraph{EnergyService}
Contiene l'intera logica di business estratta dal controller nel refactoring architetturale. È il cuore computazionale del sistema.
\begin{itemize}
    \item \textbf{Responsabilità:} Aggregazione giornaliera dei dati, analisi batch (best/worst day, anomalie), calcoli economici, analisi ambientali, forecasting, analisi AI.
    \item \textbf{Ottimizzazione:} I metodi \texttt{getBatchInsights()} e \texttt{getPeakHours()} sono annotati con \texttt{@Cacheable} (Spring Cache), poiché elaborano l'intero dataset senza parametri variabili e i risultati possono essere riusati senza ricalcolo.
\end{itemize}

\paragraph{RenewableRepository}
Interfaccia JPA che incapsula le operazioni di query al database PostgreSQL. Offre metodi per filtrare dati per range temporale e recuperare tutti i dati.
\begin{itemize}
    \item \textbf{Metodi Chiave:} \texttt{findAll()}, \texttt{findByDateTimeBetween(start, end)}.
    \item \texttt{Hibernate} genera automaticamente le query SQL ottimizzate.
\end{itemize}

\paragraph{LlmService}
Servizio dedicato all'integrazione con \textbf{Groq API}, un servizio cloud LLM ad alte prestazioni (LPU hardware). Costruisce un prompt strutturato nel formato OpenAI-compatible e gestisce autenticazione, serializzazione e fallback tra modelli.
\begin{itemize}
    \item \textbf{Metodo Principale:} \texttt{askAi(String promptData)}
    \item \textbf{Responsabilità:}
    \begin{enumerate}
        \item Arricchisce i dati grezzi con un contesto professionale come sistema-prompt.
        \item Crea una richiesta HTTPS POST verso \texttt{https://api.groq.com/openai/v1/chat/completions}.
        \item Autentica la richiesta tramite API Key letta dalla variabile d'ambiente \texttt{GROQ\_API\_KEY}.
        \item In caso di rate limit (HTTP 429), tenta automaticamente il modello successivo nell'array di fallback.
    \end{enumerate}
    \item \textbf{Modelli (in ordine di preferenza):} \texttt{llama-3.3-70b-versatile}, \texttt{llama-3.1-8b-instant}, \texttt{mixtral-8x7b-32768}.
\end{itemize}

\paragraph{UserSession}
Mantiene il contesto di sessione dell'utente corrente, incluso il filtro temporale impostato. Si trova nel package \texttt{session} (separato dal package \texttt{service}) per chiarezza architetturale.
\begin{itemize}
    \item \textbf{Stato Mantenuto:} Date di inizio/fine per il filtro, booleano per indicare se il filtro è attivo.
    \item \textbf{Scopo:} Evitare che il client debba ripetere i parametri di filtro a ogni richiesta.
    \item \textbf{Nota:} È un \textbf{bean \texttt{@SessionScope}} di Spring: ogni sessione HTTP ottiene un'istanza separata; utenti diversi non si interferiscono.
\end{itemize}

\paragraph{SecurityConfig e JwtFilter}
Configurano e applicano il meccanismo di sicurezza JWT.
\begin{itemize}
    \item \textbf{SecurityConfig:} Definisce quale URL è pubblico (\texttt{/api/auth/**}) e quale protected (\texttt{/api/energy/**}).
    \item \textbf{JwtFilter:} Intercetta ogni richiesta, estrae il token dall'header, lo valida, e popola il SecurityContext se valido.
\end{itemize}

\subsubsection{Struttura dei Package}
Il progetto segue una struttura a package orientata al dominio (package-by-feature):

\begin{table}[H]
    \centering
    \begin{tabular}{ll}
        \toprule
        \textbf{Package} & \textbf{Contenuto} \\
        \midrule
        \texttt{controller/} & \texttt{EnergyController}, \texttt{AuthController} — strato REST \\
        \texttt{service/} & \texttt{EnergyService}, \texttt{LlmService} — logica di business \\
        \texttt{session/} & \texttt{UserSession} — stato di sessione HTTP \\
        \texttt{security/} & \texttt{JwtFilter}, \texttt{JwtUtil}, \texttt{SecurityConfig} \\
        \texttt{config/} & \texttt{DataLoader} — import CSV all'avvio \\
        \texttt{dto/} & \texttt{DailyReportDTO}, \texttt{BatchInsightsDTO} \\
        \texttt{model/} & \texttt{RenewableData} — entità JPA \\
        \texttt{repository/} & \texttt{RenewableRepository} — accesso al DB \\
        \texttt{client/} & \texttt{HeliosClient} — CLI Java \\
        \bottomrule
    \end{tabular}
    \caption{Struttura dei package del progetto Helios.}
\end{table}

\subsection{Lato Client (Frontend CLI)}
L'interfaccia utente è una \textbf{Command Line Interface (CLI)} avanzata sviluppata in Java puro. Utilizza le librerie standard \texttt{java.net.http} per comunicare con il backend.
La scelta di una CLI enfatizza la natura "Enterprise" e di amministrazione remota del tool, permettendo l'esecuzione anche su server headless.

\subsubsection{Architettura Client}
Il client è organizzato attorno a un loop principale che:
\begin{enumerate}
    \item Stampa un menu di opzioni corrispondenti a diversi endpoint.
    \item Legge l'input dell'utente.
    \item Effettua la richiesta HTTP corrispondente al server backend.
    \item Formatta e visualizza la risposta JSON in modo leggibile.
\end{enumerate}

\paragraph{Gestione del Token JWT}
Il client memorizza il token JWT in una variabile statica (\texttt{jwtToken}) e lo allega automaticamente ad ogni richiesta nel header \texttt{Authorization: Bearer <token>}.

\paragraph{Flusso Tipico}
\begin{enumerate}
    \item \textbf{Login:} Utente seleziona opzione 0, inserisce credenziali. Il client invia richiesta POST a \texttt{/api/auth/login}.
    \item \textbf{Token Salvato:} Il token JWT viene estratto dalla risposta e salvato in memoria.
    \item \textbf{Analisi Dati:} Utente seleziona operazione (es. opzione 1 per report). Il client allega il token e invia la richiesta.
    \item \textbf{Visualizzazione:} La risposta JSON viene parsedata e presentata all'utente in formato leggibile.
\end{enumerate}

% -------------------------------------------------------------------
% CAPITOLO 3
% -------------------------------------------------------------------
\chapter{Design Pattern Implementati}

Per soddisfare i requisiti di ingegneria del software, sono stati implementati diversi pattern architetturali.

\section{Remote Facade}
Per ridurre la latenza di rete e il numero di chiamate ("chattiness"), il server espone API \textit{Coarse-Grained}. Il controller \texttt{EnergyController} agisce da facciata, incapsulando la complessità della logica di business.
\begin{itemize}
    \item \textbf{Esempio:} L'endpoint \texttt{/daily-report-session} non restituisce righe grezze del database, ma un report già aggregato e calcolato (data, produzione totale, temperatura media, status).
    \item \textbf{Vantaggi:} Riduce il traffico di rete, il client non deve effettuare molteplici richieste, la logica resta centralizzata sul server.
\end{itemize}

\subsection{Implementazione Tecnica}
Il client invoca un singolo endpoint e riceve un'intera lista di report giornalieri aggregati:

\begin{lstlisting}[language=Java, caption=Remote Facade: EnergyController delega a EnergyService]
// EnergyController - strato presentazione (thin layer)
@GetMapping("/daily-report-session")
public List<DailyReportDTO> getDailyReportsWithSession() {
    return energyService.getDailyReports();
}

// EnergyService - logica di business (facade reale)
public List<DailyReportDTO> getDailyReports() {
    List<RenewableData> allData = repository.findAll();

    // Filtro per range temporale (se attivo in sessione)
    if (userSession.isFilterActive()) {
        allData = allData.stream()
            .filter(d -> !d.getDateTime().toLocalDate()
                    .isBefore(userSession.getStart())
                && !d.getDateTime().toLocalDate()
                    .isAfter(userSession.getEnd()))
            .collect(Collectors.toList());
    }

    // Aggregazione per giorno e trasformazione in DTO
    return allData.stream()
        .collect(Collectors.groupingBy(
            d -> d.getDateTime().toLocalDate()))
        .entrySet().stream()
        .map(entry -> {
            double totalProd = entry.getValue().stream()
                .mapToDouble(RenewableData::getSystemProduction).sum();
            double avgTemp = entry.getValue().stream()
                .mapToDouble(RenewableData::getAirTemperature)
                .average().orElse(0.0);
            String status = totalProd > 2000 ? "HIGH_PERFORMANCE"
                : totalProd < 500 ? "LOW_GENERATION" : "NORMAL";
            return new DailyReportDTO(
                entry.getKey().toString(), totalProd, avgTemp, status);
        })
        .sorted(Comparator.comparing(DailyReportDTO::getDate))
        .collect(Collectors.toList());
}
\end{lstlisting}

Tramite un'unica richiesta \texttt{GET /daily-report-session}, il client riceve un vettore di \texttt{DailyReportDTO} completamente elaborati e pronti per la visualizzazione.

\section{Data Transfer Object (DTO)}
I dati interni del database (Entity \texttt{RenewableData}) non vengono mai esposti direttamente ai client. Sono stati creati oggetti specifici per il trasporto dati, come \texttt{DailyReportDTO} e \texttt{BatchInsightsDTO}. 

\subsection{Vantaggi del Pattern DTO}
\begin{enumerate}
    \item \textbf{Decoupling:} Il client non conosce la struttura interna dell'entità JPA, quindi modifiche allo schema del database non rompono il contratto API.
    \item \textbf{Controllo:} Possiamo selezionare quali campi esporre e quali nascondere per ragioni di sicurezza.
    \item \textbf{Trasformazione:} Il DTO può contenere campi calcolati che non esistono nel database (es. il campo "status").
    \item \textbf{Validazione:} I DTO possono includere annotazioni di validazione Spring (\texttt{@NotNull}, \texttt{@Min}, etc.).
\end{enumerate}

\subsection{Struttura dei DTO}
\begin{lstlisting}[language=Java, caption=Esempio di DTO]
@Data
@AllArgsConstructor
public class DailyReportDTO {
    private String date;          // Data del report
    private Double totalProduction; // Produzione aggregata (kWh)
    private Double avgTemperature;  // Temperatura media (°C)
    private String status;          // NORMAL | HIGH_PERFORMANCE | LOW_GENERATION
}
\end{lstlisting}

Nota i campi calcolati come \texttt{status}, che non esistono nella tabella \texttt{renewable\_data}, ma vengono generati dal server per fornire informazioni di alto livello al client.

\section{Session State Pattern}
Sebbene REST sia stateless per design, è stato implementato un meccanismo di sessione per migliorare l'esperienza utente nelle analisi continuative.

\subsection{Problema Risolto}
In un'architettura REST pura, ogni richiesta deve contenere tutti i parametri necessari. Se vogliamo analizzare i dati di un mese intero (es. gennaio 2017), dovremmo includere i parametri \texttt{start} e \texttt{end} in \textbf{ogni} richiesta:

\begin{lstlisting}[language=bash]
GET /api/energy/daily-report?start=2017-01-01&end=2017-01-31
GET /api/energy/wind-impact?start=2017-01-01&end=2017-01-31
GET /api/energy/forecast?start=2017-01-01&end=2017-01-31
\end{lstlisting}

\subsection{Soluzione: Session State}
Il sistema implementa un meccanismo di sessione che centralizza i filtri:

\begin{enumerate}
    \item \textbf{Impostazione Filtro:} Il client invia una richiesta POST a \texttt{/session/filter} con i parametri temporali:
    \begin{lstlisting}[language=bash]
POST /api/energy/session/filter?start=2017-01-01&end=2017-01-31
    \end{lstlisting}
    
    \item \textbf{Salvataggio in Sessione:} Il server estrae il range e lo memorizza in un bean \texttt{UserSession} scoped alla sessione HTTP:
    \begin{lstlisting}[language=Java]
@PostMapping("/session/filter")
public String setSessionFilter(
        @RequestParam String start, 
        @RequestParam String end) {
    userSession.setRange(
        LocalDate.parse(start), 
        LocalDate.parse(end)
    );
    return "Filtro sessione attivato: dal " 
        + start + " al " + end;
}
    \end{lstlisting}
    
    \item \textbf{Riutilizzo:} Le richieste successive recuperano automaticamente il filtro dalla sessione:
    \begin{lstlisting}[language=bash]
GET /api/energy/daily-report-session  # Utilizza il filtro salvato
GET /api/energy/wind-impact-session   # Stessa cosa
    \end{lstlisting}
\end{enumerate}

\subsection{Implementazione Tecnica}
\begin{lstlisting}[language=Java, caption=UserSession Bean]
@Component
@Scope("session")
public class UserSession {
    private LocalDate start;
    private LocalDate end;
    private boolean filterActive = false;
    
    public void setRange(LocalDate start, LocalDate end) {
        this.start = start;
        this.end = end;
        this.filterActive = true;
    }
    
    public boolean isFilterActive() {
        return filterActive;
    }
    
    // getter/setter...
}
\end{lstlisting}

L'annotazione \texttt{@Scope("session")} del framework Spring garantisce che:
- Ogni sessione HTTP riceve un'istanza separata di \texttt{UserSession}.
- Lo stato è persistito finché la sessione rimane attiva (default 30 minuti).
- Non ci sono interferenze tra utenti diversi.

\section{Request Batch}
Per ottimizzare la scoperta di informazioni critiche, l'endpoint \texttt{/batch-suggestions} esegue molteplici query in parallelo su tutto il dataset e restituisce un unico pacchetto JSON contenente le conclusioni più importanti.

\subsection{Funzionamento}
Anziché costringere il client a effettuare 4-5 richieste separate:
\begin{lstlisting}[language=bash]
GET /best-day
GET /worst-day
GET /anomalies
GET /total-production
GET /avg-efficiency
\end{lstlisting}

Un'unica richiesta \texttt{GET /batch-suggestions} elabora tutte queste analisi congiuntamente:

\begin{lstlisting}[language=Java, caption=Request Batch - Implementazione]
@GetMapping("/batch-suggestions")
public BatchInsightsDTO getBatchSuggestions() {
    List<RenewableData> allData = repository.findAll();
    
    // Trova il giorno migliore
    RenewableData best = allData.stream()
        .max(Comparator.comparing(
            RenewableData::getSystemProduction
        ))
        .orElse(null);
    
    // Trova il giorno peggiore
    RenewableData worst = allData.stream()
        .filter(d -> d.getSystemProduction() > 0)
        .min(Comparator.comparing(
            RenewableData::getSystemProduction
        ))
        .orElse(null);
    
    // Calcola produzione totale
    double totalSum = allData.stream()
        .mapToDouble(RenewableData::getSystemProduction)
        .sum();
    
    // Rileva anomalie (alta radiazione, bassa produzione)
    RenewableData anomaly = allData.stream()
        .filter(d -> d.getRadiation() > 100 
                  && d.getSystemProduction() < 50)
        .findFirst()
        .orElse(null);
    
    String anomalyMsg = anomaly != null 
        ? "ALERT: Giorno " + anomaly.getDateTime().toLocalDate()
            + " - Radiazione: " + anomaly.getRadiation() + " ma Produzione: "
            + anomaly.getSystemProduction()
        : "Nessuna anomalia rilevata.";
    
    return new BatchInsightsDTO(
        best != null ? best.getDateTime().toLocalDate() : "N/A",
        best != null ? best.getSystemProduction() : 0.0,
        worst != null ? worst.getDateTime().toLocalDate() : "N/A",
        worst != null ? worst.getSystemProduction() : 0.0,
        anomalyMsg,
        totalSum
    );
}
\end{lstlisting}

\subsection{Risposta Batch}
Una sola risposta JSON contiene tutte le informazioni critiche:
\begin{lstlisting}[language=json, caption=Risposta /batch-suggestions]
{
    "bestDay": "2017-06-21",
    "bestProduction": 3245.67,
    "worstDay": "2017-12-05",
    "worstProduction": 89.23,
    "anomalyMessage": "ALERT: Giorno 2017-03-15 - Radiazione: 650 ma Produzione: 25",
    "totalProduction": 456789.12
}
\end{lstlisting}

Questo riduce drasticamente il numero di round-trip HTTP e la latenza globale percepita dal client.

\section{Adapter Pattern - LLM Integration}

\subsection{Problema e Contesto}
L'integrazione con un'API di generazione del linguaggio naturale richiede adattare il formato dei dati energetici a quello atteso dall'API remota, incapsulando il protocollo HTTPS, l'autenticazione Bearer e il formato OpenAI-compatible in modo trasparente al resto dell'applicazione.

\subsection{Soluzione: LlmService come Adapter}
Il servizio \texttt{LlmService} funge da adapter tra il dominio energetico (\texttt{EnergyService}) e il dominio AI (\textbf{Groq API}). Il chiamante non sa nulla del provider, del modello specifico utilizzato, né del protocollo di rete.

\begin{lstlisting}[language=Java, caption=LlmService - Adapter verso Groq API con fallback automatico]
@Service
public class LlmService {

    private static final String GROQ_URL =
        "https://api.groq.com/openai/v1/chat/completions";

    // Fallback automatico: qualita' -> velocita'
    private static final String[] MODELS = {
        "llama-3.3-70b-versatile",
        "llama-3.1-8b-instant",
        "mixtral-8x7b-32768"
    };

    @Value("${groq.api.key:}")
    private String apiKey;

    public String askAi(String promptData) {
        for (String model : MODELS) {
            try {
                // Costruisce la richiesta nel formato OpenAI-compatible
                Map<String, Object> userMsg = Map.of(
                    "role", "user",
                    "content", "Sei un ingegnere energetico esperto. "
                        + "Analizza i dati e dai consigli tecnici:\n"
                        + promptData
                );
                Map<String, Object> body = Map.of(
                    "model", model,
                    "messages", List.of(userMsg),
                    "temperature", 0.7,
                    "max_tokens", 500
                );
                // Invia la richiesta HTTPS con Bearer token
                HttpHeaders headers = new HttpHeaders();
                headers.setBearerAuth(apiKey);
                headers.setContentType(MediaType.APPLICATION_JSON);
                ResponseEntity<Map> response = restTemplate.exchange(
                    GROQ_URL, HttpMethod.POST,
                    new HttpEntity<>(body, headers), Map.class
                );
                return extractContent(response.getBody());

            } catch (HttpClientErrorException e) {
                // 429 = rate limit: prova il modello successivo
                if (e.getStatusCode().value() == 429) continue;
                return "Errore API: " + e.getMessage();
            }
        }
        return "Tutti i modelli Groq hanno restituito un errore.";
    }
}
\end{lstlisting}

\subsection{Flusso di Richiesta All'LLM}

Quando il client richiede l'analisi AI tramite \texttt{/smart-analysis}:

\begin{enumerate}
    \item \textbf{Request HTTP:}
    \begin{lstlisting}[language=bash]
GET /api/energy/smart-analysis?start=2017-06-01&end=2017-06-30
Authorization: Bearer <jwt_token>
    \end{lstlisting}

    \item \textbf{Elaborazione Backend (EnergyService):} Il servizio recupera max 15 record dal repository per il periodo richiesto, li formatta in testo strutturato e li passa a \texttt{LlmService.askAi()}.

    \item \textbf{Body POST inviato a Groq (formato OpenAI-compatible):}
    \begin{lstlisting}[language=json, caption=Body POST a Groq API]
{
    "model": "llama-3.3-70b-versatile",
    "messages": [{
        "role": "user",
        "content": "Sei un ingegnere energetico esperto...\n
            - Data: 2017-06-15, Prod: 2345.67, Rad: 850.0, ..."
    }],
    "temperature": 0.7,
    "max_tokens": 500
}
    \end{lstlisting}

    \item \textbf{Risposta Groq (formato OpenAI-compatible):}
    \begin{lstlisting}[language=json, caption=Risposta da Groq API]
{
    "choices": [{
        "message": {
            "content": "Analizzando i dati, osservo una produzione
                stabile con un picco il 17 giugno. La radiazione
                elevata e le temperature moderate sono favorevoli.
                Consiglio: effettuare pulizia dei moduli."
        }
    }]
}
    \end{lstlisting}

    \item \textbf{Risposta al Client:} Il backend estrae \texttt{choices[0].message.content} e lo restituisce come plain text al browser in 1-2 secondi (Groq LPU hardware).
\end{enumerate}

\subsection{Configurazione dell'API Key}
La chiave API Groq non viene mai hardcodata nel codice sorgente. Viene letta esclusivamente dalla variabile d'ambiente \texttt{GROQ\_API\_KEY}:
\begin{lstlisting}[language=bash, caption=Configurazione sicura tramite variabile d'ambiente]
# application.properties (default vuoto: se assente, LLM disabilitata)
groq.api.key=${GROQ_API_KEY:}

# Avvio con chiave
GROQ_API_KEY=gsk_... mvn spring-boot:run
\end{lstlisting}
Modelli disponibili con fallback automatico in caso di rate limit (HTTP 429): \texttt{llama-3.3-70b-versatile} (primario, 70B parametri), \texttt{llama-3.1-8b-instant} (veloce), \texttt{mixtral-8x7b-32768} (alternativo).

\section{Remote Proxy}

\subsection{Definizione e Contesto}
Il pattern \textbf{Remote Proxy} (GoF - Structural) fornisce un surrogato locale di un oggetto che risiede su un sistema remoto. Il proxy gestisce tutta la comunicazione di rete in modo trasparente, così il codice chiamante interagisce con esso come se fosse un oggetto locale.

La distinzione con il \textbf{Remote Facade} è fondamentale:
\begin{itemize}
    \item \textbf{Remote Proxy:} rappresenta un \textit{singolo} oggetto remoto, nascondendo la complessità della rete. L'interfaccia è identica (o analoga) a quella dell'oggetto reale.
    \item \textbf{Remote Facade:} riduce il numero di round-trip aggregando \textit{più operazioni} in una sola chiamata. L'interfaccia è deliberatamente più grossa di quella degli oggetti sottostanti.
\end{itemize}

Nel progetto Helios sono presenti \textbf{due istanze} del Remote Proxy.

\subsection{Remote Proxy \#1: LlmService (Proxy del modello AI remoto)}

Il componente \texttt{EnergyController} invoca \texttt{llmService.askAi(data)} come una semplice chiamata locale. \texttt{LlmService} è il proxy che nasconde completamente ciò che avviene "dall'altra parte della rete": la connessione HTTP all'API Groq, la serializzazione JSON del body, l'autenticazione Bearer Token, la selezione del modello con fallback automatico, e la deserializzazione della risposta.

\begin{lstlisting}[language=Java, caption=LlmService come Remote Proxy del servizio AI]
@Service
public class LlmService {

    private static final String GROQ_URL =
        "https://api.groq.com/openai/v1/chat/completions";

    // Fallback automatico: qualita' -> velocita'
    private static final String[] MODELS = {
        "llama-3.3-70b-versatile",
        "llama-3.1-8b-instant",
        "mixtral-8x7b-32768"
    };

    public String askAi(String promptData) {
        // Il chiamante (EnergyController) non sa nulla di:
        // - quale modello viene usato
        // - come viene serializzata la richiesta HTTP
        // - come viene gestito il fallback tra modelli
        // - come viene autenticata la chiamata
        for (String model : MODELS) {
            // Costruisce e invia la richiesta HTTP all'API remota
            // Gestisce automaticamente errori e retry
        }
    }
}
\end{lstlisting}

Dal punto di vista del controller, la chiamata \texttt{llmService.askAi(data)} e' identica a una chiamata locale: nessun dettaglio di rete, nessun modello specifico, nessun header HTTP e' visibile.

\subsection{Remote Proxy \#2: HeliosClient (Proxy del server backend)}

Il client CLI espone all'operatore operazioni come \texttt{performLogin()}, \texttt{sendGetRequest()} e \texttt{setSessionFilter()} che appaiono come funzioni locali, ma internamente effettuano tutte richieste HTTP verso \texttt{http://localhost:8080/api}. L'utente interagisce con il client come se stesse usando direttamente il sistema; \texttt{HeliosClient} e' il proxy del server remoto.

\begin{lstlisting}[language=Java, caption=HeliosClient come Remote Proxy del server Spring Boot]
public class HeliosClient {

    private static final String BASE_URL = "http://localhost:8080/api";
    private static final HttpClient client = HttpClient.newHttpClient();

    // Proxy del login remoto: il chiamante non gestisce HTTP
    private static void performLogin(Scanner scanner) {
        HttpRequest request = HttpRequest.newBuilder()
                .uri(URI.create(BASE_URL + "/auth/login"))
                .header("Content-Type",
                    "application/x-www-form-urlencoded")
                .POST(HttpRequest.BodyPublishers.ofString(formData))
                .build();

        HttpResponse<String> response =
            client.send(request, HttpResponse.BodyHandlers.ofString());
        // Estrae il token JWT dalla risposta e lo salva localmente
    }

    // Proxy della chiamata GET remota: gestisce token e cookie
    private static void sendGetRequest(String endpoint) {
        HttpRequest.Builder builder = HttpRequest.newBuilder()
                .uri(URI.create(BASE_URL + endpoint))
                .header("Authorization", "Bearer " + jwtToken)
                .GET();
        // Aggiunge cookie di sessione se presente
        // Gestisce codici di risposta HTTP
    }
}
\end{lstlisting}

\subsection{Confronto tra le due istanze}

\begin{table}[H]
    \centering
    \begin{tabular}{lll}
        \toprule
        \textbf{Aspetto} & \textbf{LlmService} & \textbf{HeliosClient} \\
        \midrule
        Oggetto remoto rappresentato & API Groq (LLM) & Server Spring Boot \\
        Chiamante & EnergyController & Operatore (CLI) \\
        Protocollo nascosto & HTTPS + Bearer Auth & HTTP + JWT + Cookie \\
        Complessita' nascosta & Fallback multi-modello & Gestione token e sessione \\
        \bottomrule
    \end{tabular}
    \caption{Confronto tra le due istanze del Remote Proxy nel sistema Helios.}
\end{table}

\section{Spring Cache - Ottimizzazione delle Analisi Costose}

Alcuni endpoint elaborano l'intero dataset (oltre 225.000 record) senza parametri variabili. Eseguire queste analisi ad ogni richiesta è inutilmente costoso. Spring Cache permette di memorizzare il risultato alla prima invocazione e restituirlo direttamente nelle chiamate successive.

\subsection{Implementazione}
L'annotazione \texttt{@Cacheable} è applicata a livello di \texttt{EnergyService}, garantendo che la cache sia trasparente sia al controller che ai test:

\begin{lstlisting}[language=Java, caption=Spring Cache su EnergyService]
@Service
public class EnergyService {

    @Cacheable("batchInsights")
    public BatchInsightsDTO getBatchInsights() {
        // Analisi costosa sull'intero dataset:
        // best day, worst day, anomalie, totale produzione
        // Eseguita una sola volta; poi servita dalla cache
        List<RenewableData> allData = repository.findAll();
        ...
    }

    @Cacheable("peakHours")
    public String getPeakHours() {
        // Aggregazione oraria sull'intero storico
        // Il risultato non cambia finche' il dataset non varia
        ...
    }
}
\end{lstlisting}

La cache è abilitata a livello applicativo tramite \texttt{@EnableCaching} sulla classe principale \texttt{IsdProjectHeliosApplication}. Gli endpoint \texttt{/batch-suggestions} e \texttt{/peak-hours} beneficiano di questa ottimizzazione; gli endpoint parametrizzati (es. \texttt{/wind-impact?start=...}) non vengono cachati perché il risultato varia in base ai parametri.

% -------------------------------------------------------------------
% CAPITOLO 4 (AGGIORNATO)
% -------------------------------------------------------------------
\chapter{Analisi Dati e Funzionalità}

Il valore aggiunto di Helios risiede nei suoi algoritmi di analisi. Di seguito le principali metriche implementate.

\section{Analisi Economica ed Ecologica}
Il sistema traduce i kWh prodotti in valore monetario e impatto ambientale.
\begin{lstlisting}[language=Java, caption=Algoritmo di calcolo ROI e CO2]
double totalValue = totalKwh * 0.20; // Prezzo medio energia
double co2Saved = totalKwh * 0.4;    // Fattore emissione mix energetico
\end{lstlisting}
Dai dati analizzati, l'impianto ha generato un risparmio stimato di oltre \textbf{45.000 \texteuro} e ha evitato l'emissione di \textbf{90 tonnellate di $CO_2$}.

\section{Correlazioni Ambientali (Vento e Umidità)}
Il sistema analizza come le variabili meteo influenzano la produzione.
\begin{itemize}
    \item \textbf{Impatto Vento:} È stato rilevato che venti superiori a 4 m/s aumentano l'efficienza dei pannelli (raffreddamento naturale), incrementando la produzione media in modo significativo.
    \item \textbf{Impatto Umidità:} Al contrario, alti tassi di umidità sono stati correlati a una diminuzione della resa, spesso associata a foschia o nuvolosità.
\end{itemize}

\section{Intelligenza Artificiale Generativa}
Tramite l'endpoint \texttt{/smart-analysis}, il sistema invia un sottoinsieme di dati critici (max 15 record) all'API LLM cloud \textbf{Groq}.
\begin{itemize}
    \item \textbf{Input:} Dati grezzi (es. "Prod: 0, Radiazione: 600").
    \item \textbf{Output:} Report discorsivo (es. "Attenzione: rilevata anomalia critica. Irraggiamento presente ma produzione assente. Probabile guasto inverter.").
\end{itemize}

\section{Forecasting e Ore di Picco}
Il sistema offre funzionalità predittive basate su medie mobili storiche e identifica le fasce orarie di picco per permettere strategie di \textit{Load Shifting} (spostamento dei carichi industriali).

% -------------------------------------------------------------------
% CAPITOLO 5
% -------------------------------------------------------------------
\chapter{Sicurezza e Flusso Dati}

La sicurezza di Helios è progettata attorno al paradigma \textbf{Stateless}. A differenza delle sessioni server-side tradizionali (che richiedono memoria sul server per ogni utente attivo), l'architettura RESTful adotta un approccio basato su \textbf{Token JWT (JSON Web Token)}. 

Questo garantisce la scalabilità orizzontale: ogni richiesta HTTP è autosufficiente e contiene tutte le informazioni necessarie per identificare il richiedente.

\section{Fase 1: Negoziazione del Token (Handshake)}
Il flusso di sicurezza inizia quando un operatore tenta di accedere al sistema. Poiché il server non mantiene memoria delle connessioni passate, è necessaria una fase preliminare di autenticazione.

\begin{enumerate}
    \item \textbf{Invio Credenziali:} Il Client CLI raccoglie username e password e invia una richiesta \texttt{POST} all'endpoint pubblico \texttt{/api/auth/login}.
    \item \textbf{Verifica e Firma:} Il server verifica l'hash della password nel database. Se corretto, genera un token JWT firmato con algoritmo \texttt{HMAC-SHA256} contenente:
    \begin{itemize}
        \item \textit{Subject:} L'identificativo dell'utente (es. "admin").
        \item \textit{Expiration:} La data di scadenza (fissata a 24 ore).
        \item \textit{Signature:} Una firma crittografica per prevenire manomissioni.
    \end{itemize}
    \item \textbf{Consegna:} Il token viene restituito al client, che lo memorizza in una variabile locale per le richieste future.
\end{enumerate}

\section{Fase 2: Il Viaggio della Richiesta}
Una volta autenticato, ogni volta che il client richiede un report (es. \texttt{/daily-report}), deve presentare il "pass" digitale ottenuto. 

La Figura \ref{fig:high_level_flow} illustra il percorso fisico e logico che una richiesta compie all'interno dello stack tecnologico Spring Boot.

\begin{figure}[H]
    \centering
    \begin{tikzpicture}[node distance=1.5cm, auto,
        block/.style={rectangle, draw, fill=blue!5, text width=12em, text centered, rounded corners, minimum height=3em, drop shadow},
        cloud/.style={draw, ellipse, fill=red!10, node distance=3cm, minimum height=2em},
        arrow/.style={thick, ->, >=stealth}
    ]

    % Nodi
    \node [cloud] (client) {\textbf{Client CLI}};
    \node [block, below=1cm of client] (tomcat) {1. \textbf{Tomcat Container}\\(Ricezione Pacchetto HTTP)};
    \node [block, below=0.8cm of tomcat] (dispatcher) {2. \textbf{Dispatcher Servlet}\\(Routing Centrale)};
    \node [block, below=1cm of dispatcher, fill=orange!10, text width=14em] (security) {3. \textbf{Security Filter Chain}\\(Ispezione e Validazione)};
    \node [block, below=1cm of security, fill=green!10] (controller) {4. \textbf{Energy Controller}\\(Logica di Business)};

    % Archi descrittivi
    \draw [arrow] (client) -- node [right] {\footnotesize Header: \texttt{Authorization: Bearer <token>}} (tomcat);
    \draw [arrow] (tomcat) -- (dispatcher);
    \draw [arrow] (dispatcher) -- node [right] {\footnotesize Intercettazione} (security);
    \draw [arrow] (security) -- node [right] {\footnotesize Utente Autenticato} (controller);
    
    % Ritorno
    \draw [->, dashed, color=gray] (controller.east) to[out=0,in=0] node [right] {\footnotesize JSON Response} (client.east);

    \end{tikzpicture}
    \caption{Panoramica del flusso di richiesta: dal Client al Controller attraverso i layer di sicurezza.}
    \label{fig:high_level_flow}
\end{figure}

Come evidenziato nello schema:
\begin{itemize}
    \item La richiesta viene intercettata dalla \textbf{Security Chain} prima ancora che il Controller sappia della sua esistenza.
    \item Se la catena di sicurezza fallisce, la richiesta viene respinta immediatamente, risparmiando risorse di elaborazione.
\end{itemize}

\section{Fase 3: Il Filtro JWT (Deep Dive)}
Il componente critico dell'architettura è il \texttt{JwtFilter}. Questo filtro agisce come un "doganiere" all'ingresso del sistema.
La logica decisionale, illustrata nella Figura \ref{fig:security_chain_detail}, segue un algoritmo rigoroso per determinare se lasciar passare la richiesta o bloccarla.

\begin{figure}[H]
    \centering
    \resizebox{1\textwidth}{!}{ 
    \begin{tikzpicture}[node distance=2cm, auto,
        decision/.style={diamond, draw, fill=yellow!20, text width=5em, text bad centered, inner sep=0pt},
        block/.style={rectangle, draw, fill=blue!10, text width=9em, text centered, rounded corners, minimum height=3em},
        line/.style={draw, thick, -latex'},
        error/.style={rectangle, draw, fill=red!20, text centered, rounded corners, minimum height=3em}
    ]

    % --- NODI ---
    \node [block, fill=gray!20] (start) {\textbf{Ingresso Richiesta}};
    
    \node [decision, below=1cm of start] (hasheader) {Header "Bearer"\\Presente?};
    
    \node [decision, below=2cm of hasheader] (isvalid) {Firma e Data\\Valide?};
    
    \node [block, right=3cm of isvalid, fill=green!10] (context) {\textbf{Popola SecurityContext}\\(Utente Identificato)};
    
    \node [block, below=2.5cm of isvalid, fill=orange!10] (nextfilter) {\textbf{Authorization Filter}\\(Controllo Ruoli)};
    
    % Esiti negativi
    \node [error, left=2cm of hasheader] (anon) {\textbf{Utente Anonimo}\\(Nessuna identità)};
    \node [error, right=3cm of hasheader] (invalid) {\textbf{Token Corrotto}\\(Log Errore)};

    % Esito finale
    \node [block, below=1.5cm of nextfilter, fill=green!20] (success) {\textbf{Accesso Concesso}\\(Al Controller)};
    \node [error, left=2cm of nextfilter] (denied) {\textbf{403 Forbidden}\\(Blocco)};

    % --- FLUSSI ---
    \path [line] (start) -- (hasheader);
    
    % Check Header
    \path [line] (hasheader) -- node [right] {Sì} (isvalid);
    \path [line] (hasheader) -- node [above] {No} (anon);
    \path [line] (anon) |- (nextfilter); % Gli anonimi proseguono, ma verranno bloccati dopo se la risorsa è protetta
    
    % Check Token
    \path [line] (isvalid) -- node [above] {Sì} (context);
    \path [line] (isvalid) -| node [near start, above] {No (Scaduto/Falso)} (invalid);
    
    % Context to Next
    \path [line] (context) |- (nextfilter);
    
    % Final Check
    \path [line] (nextfilter) -- node [right] {Utente Autenticato?} (success);
    \path [line] (nextfilter) -- node [above] {No} (denied);

    \end{tikzpicture}
    }
    \caption{Logica decisionale interna al JwtFilter. Si noti come l'assenza del token non causi un errore immediato, ma porti al blocco nel filtro successivo (Authorization Filter) se la risorsa è protetta.}
    \label{fig:security_chain_detail}
\end{figure}

\subsection{Dettaglio delle Operazioni}
\begin{enumerate}
    \item \textbf{Estrazione:} Il filtro ispeziona l'header HTTP cercando la stringa \texttt{Authorization}.
    \item \textbf{Validazione:} Se trova un token, utilizza la chiave segreta (Secret Key) per ricalcolare la firma. Se la firma calcolata corrisponde a quella del token, significa che il dato non è stato alterato.
    \item \textbf{Contestualizzazione:} Se il token è valido, il filtro crea un oggetto \texttt{Authentication} e lo inserisce nel \texttt{SecurityContextHolder}. 
    \item \textbf{Autorizzazione:} Solo a questo punto Spring Security controlla se l'utente (ora identificato o rimasto anonimo) ha i permessi per l'URL richiesto. Se un utente anonimo tenta di accedere a \texttt{/api/energy}, riceverà un errore \textbf{403 Forbidden}.
\end{enumerate}
% -------------------------------------------------------------------
% CAPITOLO 6
% -------------------------------------------------------------------
\chapter{Implementazione Client CLI}

Il Client sviluppato (\texttt{HeliosClient.java}) funge da console di comando per l'operatore.

\section{Struttura del Menu}
Il client presenta un menu interattivo che guida l'utente attraverso le funzionalità:

\begin{lstlisting}[language=bash, caption=Esempio di Menu CLI]
MENU PRINCIPALE:
✅ LOGGATO (Token Attivo)
0. LOGIN (Richiesto!)
1. Report Semplice
2. Analisi Batch
...
7. Report Economico
9. CHIEDI ALL'IA (Groq AI)
q. Esci
\end{lstlisting}

\section{Gestione dello Stato}
Il client è stato progettato per essere "smart":
\begin{itemize}
    \item Gestisce automaticamente il ciclo di vita del \textbf{Token JWT}.
    \item Gestisce la persistenza dei \textbf{Cookie di Sessione} per i filtri temporali.
    \item Include logica di formattazione per presentare i JSON ricevuti in modo leggibile.
\end{itemize}

% -------------------------------------------------------------------
% CAPITOLO 7: TESTING, DEBUGGING E OTTIMIZZAZIONI
% -------------------------------------------------------------------
\chapter{Testing, Debugging e Ottimizzazioni}

Di seguito la validazione end-to-end del sistema e lo stack tecnologico finale adottato.

\section{Validazione End-to-End}

\subsection{Endpoints Testati}
Tutti gli endpoint sono stati testati e validati:

\begin{table}[H]
    \centering
    \begin{tabular}{lccl}
        \toprule
        \textbf{Endpoint} & \textbf{Metodo} & \textbf{Status} & \textbf{Output} \\
        \midrule
        \texttt{/api/auth/login} & POST & ✅ 200 & JWT Token \\
        \texttt{/api/energy/daily-report-session} & GET & ✅ 200 & JSON Array \\
        \texttt{/api/energy/batch-suggestions} & GET & ✅ 200 & JSON Object \\
        \texttt{/api/energy/wind-impact} & GET & ✅ 200 & Plain Text \\
        \texttt{/api/energy/forecast} & GET & ✅ 200 & Plain Text \\
        \texttt{/api/energy/financial-report} & GET & ✅ 200 & Plain Text \\
        \texttt{/api/energy/peak-hours} & GET & ✅ 200 & Plain Text \\
        \texttt{/api/energy/smart-analysis} & GET & ✅ 200 & AI Response \\
        \bottomrule
    \end{tabular}
    \caption{Stato degli endpoint API dopo ottimizzazione.}
    \label{tbl:endpoints}
\end{table}

\subsection{Screenshot delle Risposte}

\subsubsection{Login e Autenticazione}
\begin{figure}[H]
    \centering
    \fbox{
        \begin{minipage}[c][6cm][c]{0.85\textwidth}
            \centering
            \Large \textbf{Screenshot: POST /api/auth/login}
            
            \vspace{1cm}
            \small
            Mostrare il JWT token ricevuto e il messaggio di successo
            
            \vspace{1cm}
            Credenziali: \texttt{admin / password}
        \end{minipage}
    }
    \caption{Risposta all'endpoint di login con token JWT.}
    \label{fig:screenshot_login}
\end{figure}

\subsubsection{Report Giornaliero}
\begin{figure}[H]
    \centering
    \fbox{
        \begin{minipage}[c][6cm][c]{0.85\textwidth}
            \centering
            \Large \textbf{Screenshot: GET /api/energy/daily-report-session}
            
            \vspace{1cm}
            \small
            Mostrare la tabella con date, produzione totale, temperatura media e status
            
            \vspace{1cm}
            Formato: JSON Array con campi date, totalProduction, avgTemperature, status
        \end{minipage}
    }
    \caption{Report giornaliero aggregato. I dati vengono raggruppati per giorno e calcolati in backend (Remote Facade pattern).}
    \label{fig:screenshot_daily}
\end{figure}

\subsubsection{Analisi Batch}
\begin{figure}[H]
    \centering
    \fbox{
        \begin{minipage}[c][6cm][c]{0.85\textwidth}
            \centering
            \Large \textbf{Screenshot: GET /api/energy/batch-suggestions}
            
            \vspace{1cm}
            \small
            Mostrare JSON con: bestDay, bestProduction, worstDay, worstProduction, anomalyMessage, totalProduction
            
            \vspace{1cm}
            Un'unica richiesta che raccoglie più informazioni critiche (Request Batch pattern)
        \end{minipage}
    }
    \caption{Analisi batch che racchiude multiple query aggregate in una singola risposta.}
    \label{fig:screenshot_batch}
\end{figure}

\subsubsection{Analisi Vento}
\begin{figure}[H]
    \centering
    \fbox{
        \begin{minipage}[c][6cm][c]{0.85\textwidth}
            \centering
            \Large \textbf{Screenshot: GET /api/energy/wind-impact}
            
            \vspace{1cm}
            \small
            Mostrare output formattato con:
            
            ANALISI VENTO (X dati analizzati):
            
            - Prod. Media Vento Forte: X.XX kWh
            
            - Prod. Media Vento Debole: X.XX kWh
            
            Differenza di resa: X.X\%
        \end{minipage}
    }
    \caption{Analisi dell'impatto del vento sulla produzione energetica.}
    \label{fig:screenshot_wind}
\end{figure}

\subsubsection{Report Economico}
\begin{figure}[H]
    \centering
    \fbox{
        \begin{minipage}[c][6cm][c]{0.85\textwidth}
            \centering
            \Large \textbf{Screenshot: GET /api/energy/financial-report}
            
            \vspace{1cm}
            \small
            Mostrare output formattato con:
            
            REPORT ECONOMICO \& ECOLOGICO:
            
            - Energia Totale Prodotta: X.XX kWh
            
            - Valore Economico Generato: X.XX € (a 0.20 €/kWh)
            
            - CO2 Risparmiata: X.XX kg
        \end{minipage}
    }
    \caption{Report economico ed ecologico (bug di formattazione risolto).}
    \label{fig:screenshot_financial}
\end{figure}

\subsubsection{Ore di Picco}
\begin{figure}[H]
    \centering
    \fbox{
        \begin{minipage}[c][6cm][c]{0.85\textwidth}
            \centering
            \Large \textbf{Screenshot: GET /api/energy/peak-hours}
            
            \vspace{1cm}
            \small
            Mostrare output con:
            
            ANALISI ORARIA:
            
            - Ora di Picco Assoluta: HH:00
            
            - Produzione Media in quell'ora: X.XX kWh
            
            Consiglio: Programmare i macchinari intorno alle HH:00
        \end{minipage}
    }
    \caption{Identificazione delle ore di picco di produzione.}
    \label{fig:screenshot_peak}
\end{figure}

\subsubsection{Previsione di Produzione}
\begin{figure}[H]
    \centering
    \fbox{
        \begin{minipage}[c][6cm][c]{0.85\textwidth}
            \centering
            \Large \textbf{Screenshot: GET /api/energy/forecast}
            
            \vspace{1cm}
            \small
            Mostrare output con:
            
            SIMULAZIONE PREVISIONE:
            
            Basandosi sul periodo selezionato (X record utili),
            
            la produzione giornaliera attesa è: X.XX kWh
        \end{minipage}
    }
    \caption{Previsione della produzione futura basata su medie storiche.}
    \label{fig:screenshot_forecast}
\end{figure}

\subsubsection{Analisi IA (Smart Analysis)}
\begin{figure}[H]
    \centering
    \fbox{
        \begin{minipage}[c][7cm][c]{0.85\textwidth}
            \centering
            \Large \textbf{Screenshot: GET /api/energy/smart-analysis}
            
            \vspace{0.5cm}
            \small
            Mostrare l'output testuale dell'LLM (neural-chat) con analisi semantica
            
            \vspace{0.5cm}
            Esempio di risposta attesa:
            
            \small
            "Analizzando i dati forniti, osservo una produzione stabile con variabilità meteo...
            La radiazione è elevata ma il vento moderato ha raffreddato i moduli in modo benefico.
            Suggerimento: monitorare l'umidità nelle giornate nuvolose..."
            
            \vspace{0.5cm}
            Nota: Prima richiesta richiede 30-60 secondi (Ollama + neural-chat su CPU)
        \end{minipage}
    }
    \caption{Analisi intelligente tramite API LLM cloud (Groq + llama-3.3-70b-versatile). Integrazione dell'Adapter Pattern per IA generativa; risposta in 1-2 secondi.}
    \label{fig:screenshot_ai}
\end{figure}

\section{Stack Tecnologico Finale}

Dopo le ottimizzazioni, il sistema è strutturato come segue:

\begin{table}[H]
    \centering
    \begin{tabular}{ll}
        \toprule
        \textbf{Componente} & \textbf{Tecnologia} \\
        \midrule
        Backend Server & Java 21 + Spring Boot 3.5.8 \\
        Database (Runtime) & PostgreSQL 16 (Docker) \\
        Database (Test) & H2 in-memory (solo test unitari) \\
        API Style & RESTful con JWT Authentication \\
        API Documentation & Springdoc OpenAPI / Swagger UI \\
        Frontend & JavaScript SPA + Chart.js 4.4.4 \\
        LLM Integration & Groq API (llama-3.3-70b + fallback) \\
        Caching & Spring Cache (\texttt{@Cacheable}) \\
        Container Orchestration & Docker Compose \\
        CI/CD Validazione & Maven Test Suite \\
        \bottomrule
    \end{tabular}
    \caption{Stack tecnologico del sistema Helios dopo ottimizzazione.}
    \label{tbl:stack}
\end{table}

% -------------------------------------------------------------------
% CONCLUSIONI
% -------------------------------------------------------------------
\chapter{Conclusioni}

Il progetto Helios ha dimostrato con successo come un'architettura a microservizi (o client-server disaccoppiata) possa risolvere problemi reali nel dominio dell'energia.

L'integrazione di pattern consolidati come Remote Facade, Remote Proxy, DTO e Session State ha garantito performance e manutenibilità. Il refactoring architetturale (estrazione di \texttt{EnergyService}, separazione dei package per dominio) ha nettamente migliorato la leggibilità e testabilità del codice. L'adozione di tecnologie emergenti come le LLM cloud (\textbf{Groq API}) ha aggiunto un livello di intelligenza ai dati non comune nei sistemi SCADA tradizionali, con tempi di risposta di 1-2 secondi grazie all'hardware LPU.

I risultati ottenuti sul dataset di test confermano la validità economica e tecnica della soluzione per contesti industriali. La fase di sviluppo ha evidenziato l'importanza della sicurezza (gestione delle API key tramite variabili d'ambiente, JWT stateless), della robustezza (fallback multi-modello, error handling) e dell'ottimizzazione (Spring Cache per analisi costose, paginazione client-side con Chart.js) nella costruzione di sistemi distribuiti affidabili.

\end{document}

